\chapter*{Installation instructions}
\addcontentsline{toc}{chapter}{Installation instructions}

\newcommand{\GMSHver}{2.2.0}
\newcommand{\instpath}{\texttt{INSTALLPATH}}
\newcommand{\version}{\texttt{VERSION}}

In the following, \version{} denotes the version number of the \TiberCAD{} release
you downloaded and \instpath{} denotes the directory where \TiberCAD{} gets installed.
Version \GMSHver{} of GMSH (\url{http://www.geuz.org/gmsh}) will be installed together with
\TiberCAD. For the Linux version of GMSH you need OpenGL libraries installed on your system.

\section*{Prerequisites}

Get the installer package for your OS/architecture from \url{http://www.tibercad.org}
or by contacting \mailto{support@tibercad.org}.
Table~\ref{install:table:packages} lists the packages available for download.
To run \TiberCAD{} you will also need a license file that you will have to copy into
the installation directory of \TiberCAD.
\begin{table}
\begin{tabular}{l|p{8cm}}
\emph{installer package name} & \emph{Target architecture} \\
\hline  \\
\verb+tibercad_installer.exe+ & Windows 32-bit \\
\verb+tibercad_+\textsl{version}\verb+_i386.deb+ & Debian GNU/Linux 4.0 ``etch'', Intel x86 \\
\verb+tibercad-+\textsl{version}\verb+.tar.gz+ & Generic Linux package
\end{tabular}
\caption{Installer packages}
\label{install:table:packages}
\end{table}

In the Windows version, some graphical features such as graphical convergence monitors are only
available if an X Window server is installed and running.




\section*{Windows installation procedure}

To install \TiberCAD{} in Windows, run the setup program \verb+tibercad_setup.exe+.
During the installation you can choose the installation directory. After finishing
installation, copy your license file (tibercad.lic) into the \verb+license+ subdirectory
of the \TiberCAD{} installation directory (\instpath\verb+/license+), without changing
its filename.


\section*{Linux installation procedure}

The installation procedure for the Linux version of \TiberCAD{} depends on your
distribution. Download the installer package that best fits your setup.

 The standard method to launch \TiberCAD{} is by means of a shell script that is
installed alongside the TiberCAD executable. It takes care of setting all
necessary environment variables. If for some reason you have to run the executable
directly, remember to set \verb+TIBERCADROOT+ to the TiberCAD installation directory
(\instpath) and \verb+LD_LIBRARY_PATH+ to \instpath\verb+/lib+.

\subsection*{Debian}

If you are running Debian GNU/Linux 4.0 ``etch'' on a i386 architecture, you can use the
Debian package \verb+tibercad_+\textsl{version}\verb+_i386.deb+.
Install it as root using dpkg or similar:
\begin{verbatim}
# dpkg --install tibercad_VERSION_i386.deb
\end{verbatim}

  The package will be installed in /usr/share/tibercad. Next, copy your
license file (tibercad.lic) into /usr/share/tibercad/license/ without
changing the filename.


\begin{description}
\item[NOTE:] The debian version of TiberCAD depends on the following Debian packages:
 \begin{itemize}
   \item libboost-regex1.33.1
   \item libboost-filesystem1.33.1
   \item libblas.so.3 (provided e.g. by atlas3-base)
   \item liblapack.so.3 (provided e.g. by atlas3-base)
 \end{itemize}
\end{description}

\subsection*{Other Linux distributions}

If you have a distribution other than Debian 4.0 ``etch'' or you want to install \TiberCAD{}
into a different directory, then use the .tgz or .tbz installation packages.
Unpack the archive, cd to the unpacked directory tibercad-\version{} and run the
install script.

  After installation, copy your license file (tibercad.lic) into the `license'
subdirectory of the \TiberCAD{} installation directory (\instpath/license) without
changing the filename.


  

\section*{Quick start guide}

In the `examples' subdirectory you can find several examples ready to run.
They are the same as
the tutorials on \url{http://www.tibercad.org/documentation/tutorial/list}.

\subsection*{Windows}

Open Windows Explorer and go to the \TiberCAD{} installation directory. 
If you have write permission in the installation directory, you can browse to
an examples directory and start the simulation by double clicking the input file, e.g. bulk.tib in the bulk\_Si example.
If not, copy the whole directory to a location in your personal area and run the
examples from there.

  If you cannot run \TiberCAD{} by double clicking an input file (*.tib), then the
input files are probably not correctly associated with the \TiberCAD{} executable.
In this case, try to establish the association by right-clicking the input file,
choosing \verb+open with...+ $\rightarrow$ \verb+Choose Program...+ $\rightarrow$ \verb+Browse...+, browsing to the \TiberCAD{} installation directory and choosing the \TiberCAD{} executable, \verb+tibercad.exe+.
A directory containing the simulation results will be created with the name provided in the input file.


\subsection*{Linux}

After the correct installation of \TiberCAD{} you should be able to run \TiberCAD{} from the command line using the command \verb+tibercad+.
If not, you probably have to add the \verb+bin+ subdirectory of the \TiberCAD{} installation directory to your \verb+PATH+ environment variable or start the \TiberCAD{} executable using the absolute path (\instpath\verb+/bin/tibercad+).
Copy the directory of the example you want to run, e.g. bulk\_Si, to your home
directory or any place you have write permissions for. Change to the newly
created directory and run \TiberCAD{} by (assuming the bulk\_Si example)
\begin{verbatim}
$ tibercad bulk.tib
\end{verbatim}
A directory containing the simulation results will be created with the name provided in the input file.


\section*{Bug reports / Feedback}

Please send bug reports, feedback or suggestions to \mailto{support@tibercad.org}.
When submitting bug reports, please always include the full version number of \TiberCAD{} you are running.
The full version number appears in the first line of output when running the program:
\begin{verbatim}
$ tibercad

TiberCAD version 0.1.0.961

Usage: tibercad <inputfile>

\end{verbatim}
